\documentclass[10pt,letterpaper]{article}

\usepackage[letterpaper,top=0.48in,bottom=0.44in,left=0.66in,right=0.66in]{geometry}
\usepackage[T1]{fontenc}
\usepackage[utf8]{inputenc}
\usepackage{lmodern}
\usepackage{microtype}
\usepackage{xcolor}
\usepackage{enumitem}
\usepackage{needspace}
\usepackage{xurl}
\usepackage{hyperref}
\input{glyphtounicode}
\pdfgentounicode=1

\DisableLigatures{encoding = *, family = * }
\DeclareUnicodeCharacter{2014}{\textemdash}
\DeclareUnicodeCharacter{2013}{\textendash}
\DeclareUnicodeCharacter{00B7}{\textperiodcentered}
\DeclareUnicodeCharacter{201C}{\textquotedblleft}
\DeclareUnicodeCharacter{201D}{\textquotedblright}
\DeclareUnicodeCharacter{2019}{\textquoteright}
\definecolor{ink}{HTML}{18222C}
\definecolor{accent}{HTML}{185A73}
\definecolor{muted}{HTML}{485762}
\hypersetup{
  colorlinks=true,
  urlcolor=accent,
  linkcolor=accent,
  pdfauthor={Fehmi Özbayrak},
  pdftitle={Fehmi Özbayrak — Curriculum Vitae},
  pdfsubject={Curriculum Vitae},
  pdfkeywords={energy systems, reservoir engineering, geomechanics, scientific software, machine learning}
}

\pagestyle{empty}
\setlength{\parindent}{0pt}
\setlength{\parskip}{0pt}
\setlength{\emergencystretch}{2em}
\linespread{0.95}
\color{ink}
\raggedright

\newcommand{\cvsection}[1]{%
  \needspace{4\baselineskip}%
  \vspace{5pt}%
  {\fontsize{11.5}{13.2}\selectfont\bfseries\color{accent} #1}\par
  \vspace{1.2pt}\color{accent}\hrule height 0.55pt\color{ink}\vspace{3pt}%
}

\newcommand{\cventry}[4]{%
  \needspace{4\baselineskip}%
  {\bfseries #1} \textemdash{} #2\par
  {\color{muted}\itshape #3\if\relax\detokenize{#4}\relax\else{} · #4\fi}\par
  \vspace{1pt}%
}

\newcommand{\cvproject}[2]{%
  \needspace{3\baselineskip}%
  {\bfseries #1}\hspace{0.35em}{\color{muted}#2}\par
  \vspace{0.5pt}%
}

\newenvironment{cvitems}{%
  \begin{itemize}[leftmargin=13pt,label=\textbullet,itemsep=1.2pt,topsep=1pt,parsep=0pt,partopsep=0pt]
}{\end{itemize}\vspace{1.5pt}}

\newenvironment{compactitems}{%
  \begin{itemize}[leftmargin=13pt,label=\textbullet,itemsep=0.8pt,topsep=1pt,parsep=0pt,partopsep=0pt]
}{\end{itemize}\vspace{1pt}}

\begin{document}

\begin{center}
  {\fontsize{22}{25}\selectfont\bfseries Fehmi Özbayrak}\par
  \vspace{2pt}
  {\fontsize{11.5}{13}\selectfont\color{accent} Curriculum Vitae · Ph.D. Student in Petroleum and Geosystems Engineering}\par
  \vspace{4pt}
  Austin, Texas · \href{tel:+15129208349}{+1 512-920-8349} ·
  \href{mailto:ozbayrakfehmi@gmail.com}{ozbayrakfehmi@gmail.com}\par
  \href{https://www.linkedin.com/in/fozbayrak}{LinkedIn} ·
  \href{https://github.com/fozba}{GitHub} ·
  \href{https://scholar.google.com/citations?user=NdarxTgAAAAJ}{Google Scholar} ·
  \href{https://fozba.github.io/portfolio}{Portfolio}
\end{center}

\cvsection{Research Profile}
Computational energy researcher connecting reservoir engineering and geomechanics, scientific software, spatial machine learning, uncertainty quantification, and techno-economic analysis. Research spans produced-water disposal and induced seismicity, geological poromechanics, geothermal development, and data-center energy systems.

\cvsection{Education}
\cventry{Ph.D. Student, Petroleum and Geosystems Engineering}{The University of Texas at Austin}{Austin, Texas}{Expected Aug 2027}
Advisor: Dr. Mukul M. Sharma. Tentative dissertation: \emph{Multiscale Geomechanics of Produced-Water Disposal in the Delaware Basin: Fracture Containment, Ground Deformation, and Induced Seismicity}. Cumulative UT M.S./Ph.D. GPA: 3.61/4.00.

\vspace{3pt}
\cventry{MSc, Petroleum and Geosystems Engineering}{The University of Texas at Austin}{Austin, Texas}{May 2024}
Thesis: \emph{A novel spatial bagging algorithm for predictive ensemble machine learning}. Co-advisors: Dr. John T. Foster and Dr. Michael J. Pyrcz.

\vspace{3pt}
\cventry{BEng, Petroleum and Natural Gas Engineering}{Middle East Technical University}{Ankara, Türkiye}{Feb 2022}
GPA: 3.21/4.00.

\cvsection{Research and Professional Experience}
\cventry{Graduate Research Assistant}{UT Austin / Center for Integrated Seismicity Research}{Austin, Texas}{Summer 2024–Present}
\begin{cvitems}
  \item Build and calibrate coupled reservoir-geomechanics models using well logs and injection and pressure histories to investigate fracture containment, pressure transmission, surface deformation, and induced-seismicity risk from produced-water disposal.
  \item Design sensitivity studies that separate stress, operating, pressure, and leakoff controls while communicating calibration non-uniqueness and model limitations.
  \item Develop auditable geological-poromechanics workflows and verification cases for MOOSE/PorousFlow simulation.
\end{cvitems}

\cventry{Reservoir Engineering Intern}{Fervo Energy}{Houston, Texas}{May–Aug 2025}
\begin{cvitems}
  \item Architected and implemented a modular Python platform for geothermal development planning, connecting configurable well schedules and surface-facility allocation to daily power forecasts, costs, and contract economics.
  \item Integrated a machine-learning proxy for production-temperature evolution to accelerate repeated schedule evaluation.
  \item Formulated drilling start times as Pyomo optimization variables and maximized project revenue subject to power-delivery constraints; built interactive Plotly and Dash analysis tools.
\end{cvitems}

\cventry{Graduate Research Assistant}{DiReCT Consortium, The University of Texas at Austin}{Austin, Texas}{2022–2024}
\begin{cvitems}
  \item Developed Spatial Bagging, an ensemble-learning method that uses geostatistical effective sample size to incorporate spatial dependence into bootstrap sampling.
  \item Led software, experiments, and first-author writing for two peer-reviewed papers on spatial prediction and uncertainty quantification; coauthored a later Spatial Random Forest extension.
\end{cvitems}

\cventry{Drilling Engineer Intern}{Middle East Air Drilling Services}{Alaşehir, Türkiye}{2021}
\begin{cvitems}
  \item Completed geothermal drilling field training around a ZJ-40 double mobile rig, performing mud and cement calculations and formation-evaluation exercises while documenting casing, cementing, MWD, and coiled-tubing operations.
\end{cvitems}

\cventry{Undergraduate Researcher}{Middle East Technical University}{Ankara, Türkiye}{Poster presented Dec 2021}
\begin{cvitems}
  \item Created synthetic permanent-downhole-gauge data, introduced noise and outliers, and compared nonlinear regression and wavelet-based denoising methods; presented the work at the ADIMODTÜ Project Poster Presentation.
\end{cvitems}

\newpage

\cvsection{Research Software and Selected Research Projects}
\cvproject{JAVELINA — Auditable Geological Poromechanics Workflow}{2026}
\begin{cvitems}
  \item Developed a declarative Python workflow that converts geological surfaces, meshes, material data, and well schedules into auditable MOOSE/PorousFlow cases for serial and MPI execution.
  \item Implemented completion intersection, conservative signed-source allocation, provenance tracking, and automated result-quality reporting.
  \item Verified generated operators against six analytical, semi-analytical, and manufactured benchmarks, with near-second-order spatial convergence and 1.196\% maximum peak-normalized error. The project remains unpublished.
\end{cvitems}

\cvproject{Data Center Integrated Systems Model (DC-ISM)}{2025–2026}
\begin{cvitems}
  \item Built a modular, configuration-driven Python demand-side engine that maps AI workload states and hardware fleets to hourly IT, cooling, electrical, water, PUE, cost, carbon, and flexibility profiles.
  \item Implemented measured-weather ingestion, tariffs, sensitivity screening, correlated uncertainty propagation, scenario-year bootstrapping, supply-land screening, and risk-aware design comparison.
  \item Defined typed interfaces, explicit scope boundaries, a limitations register, and a 155-test suite to keep assumptions and unsupported effects auditable. The project remains private.
\end{cvitems}

\cvproject{Spatial Bagging for Prediction and Uncertainty Quantification}{2022–2025}
\begin{cvitems}
  \item Developed an ensemble-learning algorithm that adjusts bootstrap sample size using geostatistical effective sample size so spatial dependence is represented rather than treated as independent data.
  \item Evaluated prediction and uncertainty behavior across synthetic two-dimensional datasets spanning spatial-correlation and noise regimes; released the \href{https://github.com/fozba/spatialbootstrap}{public Python implementation} and published two first-author papers.
\end{cvitems}

\cvproject{Delaware Mountain Group Fracture-Containment Study}{2024–Present}
\begin{cvitems}
  \item Built and calibrated a coupled flow-and-fracture model against multi-year injection history and daily pressure observations from Delaware Basin disposal wells.
  \item Designed sensitivity studies across stress stratigraphy, injection rate and intermittency, pressure, and leakoff to distinguish controls on vertical containment, pressure buildup, and lateral fracture growth. Manuscript in preparation.
\end{cvitems}

\cvproject{Sierra Solutions — Gas-to-Data-Center Investment Model}{2026}
\begin{cvitems}
  \item Built a 15-year physics-and-finance model for a 15 MW firm / 20.25 MW installed behind-the-meter natural-gas power concept serving an AI data center, integrating redundancy, fuel logistics, contracts, debt, and project economics.
  \item Evaluated 10,000 Monte Carlo scenarios across 12 risk drivers and presented a base case with 14.47\% levered IRR, \$2.32 million NPV at 12\%, and 1.72x minimum DSCR; earned second place in the North American SPE Energython as a solo entrant.
\end{cvitems}

\cvsection{Technical Competition Projects}
\cvproject{Utah FORGE — Engineering Analysis, GEODE Engineering Datathon}{2nd place · 2025 · Two-person team}
\begin{compactitems}
  \item Wrote the analysis code integrating tracer, distributed fiber-optic temperature, strain and acoustic data, microseismicity, well surveys, and stimulation records to infer stage-level connectivity and preferential flow paths.
  \item Proposed a targeted restimulation program using perturbed-stress interpretation, oriented perforations, isolation and diversion, and pulsed initiation; recommendations were not represented as implemented results.
\end{compactitems}

\cvproject{Nevada Blind-Geothermal Exploration — GEODE Geosciences Datathon}{2nd place · 2025 · Two-person team}
\begin{compactitems}
  \item Built a spatial CNN workflow using strain, gravity, permeability, modeled temperature, earthquake-density, and fault-distance rasters to reproduce and interrogate play-fairway favorability patterns.
  \item Used geographically separated development and validation regions to identify candidate zones for follow-up without presenting them as physically validated discoveries.
\end{compactitems}

\cvproject{Hydraulic-Fracturing Fuel Forecast — UT PGE Energy AI Hackathon}{1st place · 2025 · Five-person team}
\begin{compactitems}
  \item Served as the main machine-learning architect and built most of a random-forest regression workflow forecasting grid, diesel, and CNG requirements from operational, fleet, fuel, formation, and weather features.
  \item Compared 12 CatBoost-derived features with 1,082 one-hot features and calibrated ensemble uncertainty across 100 repeated runs.
\end{compactitems}

\cvproject{Missing-Well Production Forecasting — UT Energy AI Hackathon}{Best Presentation · 2026 · Team Chain}
\begin{compactitems}
  \item Served as the main machine-learning engineer for a field-scale workflow predicting oil, water, and gas production for wells omitted from the supplied dataset.
\end{compactitems}

\cvproject{Utah FORGE Microseismic Analysis — SPE Geothermal Datathon}{2nd place · 2023 · Four-person team}
\begin{compactitems}
  \item Applied bootstrap ensemble clustering to microseismic events, clustering each resampled realization and aggregating the results to interpret event structure.
\end{compactitems}

\cvsection{Peer-Reviewed Publications}
\begin{enumerate}[leftmargin=18pt,label={\arabic*.},itemsep=3pt,topsep=1pt,parsep=0pt]
  \item \textbf{Özbayrak, F.}, Foster, J. T., and Pyrcz, M. J. (2025). ``Spatial bagging for predictive machine learning uncertainty quantification.'' \emph{Computers \& Geosciences}, 203, 105947. \href{https://doi.org/10.1016/j.cageo.2025.105947}{doi:10.1016/j.cageo.2025.105947}. First and corresponding author.
  \item Merzoug, A., \textbf{Özbayrak, F.}, Foster, J. T., and Pyrcz, M. J. (2025). ``Beyond random forest: how spatial bagging and spatial random forest dominate for subsurface applications?'' \emph{Computational Geosciences}, 29, article 52. \href{https://doi.org/10.1007/s10596-025-10388-0}{doi:10.1007/s10596-025-10388-0}.
  \item \textbf{Özbayrak, F.}, Foster, J. T., and Pyrcz, M. J. (2024). ``Spatial bagging to integrate spatial correlation into ensemble machine learning.'' \emph{Computers \& Geosciences}, 186, 105558. \href{https://doi.org/10.1016/j.cageo.2024.105558}{doi:10.1016/j.cageo.2024.105558}. First and corresponding author.
\end{enumerate}

\cvsection{Manuscripts in Preparation}
\begin{compactitems}
  \item \textbf{Özbayrak, F.} ``JAVELINA: an accessible, quality-transparent workflow for geological poromechanics with MOOSE.'' Manuscript in preparation; not submitted.
  \item \textbf{Özbayrak, F.} et al. ``Induced fracture geometry and vertical containment during shallow wastewater injection in the Delaware Mountain Group: A calibrated geomechanical modeling study.'' Manuscript in preparation; title and author order subject to coauthor review.
\end{compactitems}

\cvsection{Awards and Honors}
\begin{compactitems}
  \item Second Place, North America SPE Asset Management Technical Section Energython — Sierra Solutions, solo entry, 2026.
  \item Judges' Selection Award for Best Presentation, UT Energy AI Hackathon — Team Chain, 2026.
  \item First Place, UT PGE Energy AI Hackathon — Data Comrades, five-person team, 2025.
  \item Second Place, GEODE Engineering Datathon — Operation Hotspot, two-person team, 2025.
  \item Second Place, GEODE Geosciences Datathon — Operation Hotspot, two-person team, 2025.
  \item Second Place, SPE Geothermal Datathon — Team NDT, four-person team, 2023.
\end{compactitems}

\cvsection{Selected Technical Presentations}
\begin{compactitems}
  \item ``Sierra Solutions,'' North America SPE Asset Management Technical Section Energython Final, Jan 2026.
  \item Utah FORGE engineering analysis and Nevada blind-geothermal exploration, GEODE Datathon, 2025.
  \item Hydraulic-fracturing fuel forecasting, UT PGE Energy AI Hackathon, 2025.
  \item \textbf{Özbayrak, F.} and Gökalp Yavuz, F. ``Denoising and Outlier Removal from Permanent Downhole Gauge Using Nonlinear Regression Methods,'' ADIMODTÜ Project Poster Presentation, 2021. \href{https://hdl.handle.net/11511/97232}{hdl:11511/97232}.
\end{compactitems}

\cvsection{Teaching}
\cventry{Teaching Assistant, PGE 383/379 — High Performance Computing for Engineers}{The University of Texas at Austin}{Austin, Texas}{Spring 2024}
Supported Julia programming laboratories and occasional office hours for approximately 30 students in Dr. John T. Foster's course.

\cvsection{Leadership and Service}
\cventry{President}{Switch Energy Alliance UT Chapter}{Austin, Texas}{2026–Present}
Lead an eight-officer student organization and secured an approximately \$3,000 budget while building undergraduate representation and cross-campus energy partnerships.

\vspace{3pt}
\cventry{Vice President for External Affairs}{Switch Energy Alliance UT Chapter}{Austin, Texas}{2024–2025}

\vspace{3pt}
\cventry{Challenge Co-Architect}{Fourth Annual UT PGE Energy AI Hackathon}{Austin, Texas}{2024 event cycle}
Co-architected the challenge with Elnara Rustamzade and supported teams and event logistics for approximately 160 students across 25 teams.

\cvsection{Technical Skills}
\textbf{Programming and software:} Python, C/C++, Julia, NumPy, Pandas, scikit-learn, Pyomo, Plotly, Dash, Git, Linux/shell, MPI, LaTeX.\par
\vspace{2pt}
\textbf{Subsurface and computational methods:} Reservoir simulation, coupled flow-and-fracture modeling, finite-element poromechanics, MOOSE/PorousFlow, ResFrac, history matching, calibration, numerical verification, sensitivity analysis.\par
\vspace{2pt}
\textbf{Data and uncertainty:} Geostatistics, spatial statistics, ensemble machine learning, random forests, raster/GIS workflows, uncertainty quantification, Monte Carlo analysis, tracer/fiber-optic/microseismic interpretation.\par
\vspace{2pt}
\textbf{Energy systems:} Geothermal systems, techno-economic modeling, project finance, mathematical optimization, hourly facility-demand modeling, data-center cooling, water, PUE, tariffs, and siting analysis.

\end{document}
